\documentclass[12pt]{article}
\usepackage[utf8]{inputenc}
\usepackage[T1]{fontenc}
\usepackage{hyperref}
\title{Tugutil Language Resources for K-12 Curriculum Development}
\author{Compiled by Claude Sonnet 5 (claude-sonnet-5) for Andr\'as Kornai}
\date{}
\begin{document}
\maketitle
\section{Introduction}
Tugutil (ISO 639-3 tuj) is spoken by Togutil (O'Hongana Manyawa) communities in Halmahera, North Maluku, Indonesia. It is closely related to Tobelo. One source gives about 2,500 speakers, from a 2008 dataset; this is unverified. Care is needed not to confuse it with Tobelo (tlb). The word ``Togutil'' is considered pejorative by some of the people, who prefer O'Hongana Manyawa.
\section{Primary Educational Resources}
None found.
\section{Linguistic and Reference Materials}
\begin{itemize}
\item Wikipedia, Togutil people: \url{https://en.wikipedia.org/wiki/Togutil_people}
\item Wikipedia, Tobelo language (related language): \url{https://en.wikipedia.org/wiki/Tobelo_language}
\item Numerals data page for Tugutil: \url{https://lingweb.eva.mpg.de/channumerals/Tugutil.htm}
\end{itemize}
\section{Religious and Community Resources}
Scripture Earth has a page for code tuj: \url{https://www.scriptureearth.org/00eng.php?iso=tuj}. I could not tell from the search snippet what materials it holds.
\section{Digital and Audio Resources}
None found.
\section{Curriculum Development Context}
The language is not known to be taught in schools. Very little documentation exists, and the community faces pressure from mining on its forest territory. Any curriculum should be developed only with the community's participation and consent.
\end{document}